\qquad As Tabelas \ref{tab:comparacao_crescente}, \ref{tab:comparacao_decrescente}
e \ref{tab:comparacao_random} apresentam, lado a lado, os tempos de
execução (em segundos) obtidos pelos cinco algoritmos estudados neste
trabalho, para os seis tamanhos de instância exigidos, uma tabela para
cada um dos três tipos de entrada testados (Crescente, Decrescente e
Random). Todos os valores são os registrados nos arquivos de tempo
gerados pelo programa para cada algoritmo.

\begin{table}[htbp]
    \centering
    \begin{tabular}{|l|r|r|r|r|r|r|}
        \hline
         & 10 & 100 & 1.000 & 10.000 & 100.000 & 1.000.000 \\
        \hline
        Insertion Sort & 0,000010 & 0,000001 & 0,000004 & 0,000015 & 0,000074 & 0,000584 \\
        \hline
        Selection Sort & 0,000002 & 0,000007 & 0,000896 & 0,093473 & 9,351039 & 934,897502 \\
        \hline
        Bubble Sort & 0,000001 & 0,000001 & 0,000000 & 0,000003 & 0,000033 & 0,000321 \\
        \hline
        Shell Sort & 0,000003 & 0,000001 & 0,000005 & 0,000065 & 0,000690 & 0,008251 \\
        \hline
        Merge Sort & 0,000002 & 0,000002 & 0,000015 & 0,000194 & 0,002070 & 0,021293 \\
        \hline
    \end{tabular}
    \caption{Tempo de execução (em segundos) dos 5 algoritmos -- entrada Crescente}
    \label{tab:comparacao_crescente}
\end{table}


\begin{table}[htbp]
    \centering
    \begin{tabular}{|l|r|r|r|r|r|r|}
        \hline
         & 10 & 100 & 1.000 & 10.000 & 100.000 & 1.000.000 \\
        \hline
        Insertion Sort & 0,000007 & 0,000006 & 0,000174 & 0,015947 & 1,620129 & 166,979744 \\
        \hline
        Selection Sort & 0,000002 & 0,000008 & 0,000893 & 0,092843 & 9,336395 & 937,348384 \\
        \hline
        Bubble Sort & 0,000002 & 0,000002 & 0,000166 & 0,015925 & 1,583171 & 155,625005 \\
        \hline
        Shell Sort & 0,000002 & 0,000002 & 0,000009 & 0,000083 & 0,001082 & 0,011790 \\
        \hline
        Merge Sort & 0,000002 & 0,000002 & 0,000014 & 0,000175 & 0,001809 & 0,020753 \\
        \hline
    \end{tabular}
    \caption{Tempo de execução (em segundos) dos 5 algoritmos -- entrada Decrescente}
    \label{tab:comparacao_decrescente}
\end{table}


\begin{table}[htbp]
    \centering
    \begin{tabular}{|l|r|r|r|r|r|r|}
        \hline
         & 10 & 100 & 1.000 & 10.000 & 100.000 & 1.000.000 \\
        \hline
        Insertion Sort & 0,000002 & 0,000002 & 0,000096 & 0,008208 & 0,844211 & 80,586551 \\
        \hline
        Selection Sort & 0,000003 & 0,000010 & 0,000903 & 0,092938 & 9,362431 & 932,549777 \\
        \hline
        Bubble Sort & 0,000003 & 0,000008 & 0,000326 & 0,027070 & 8,899341 & 938,659248 \\
        \hline
        Shell Sort & 0,000003 & 0,000004 & 0,000052 & 0,000655 & 0,008226 & 0,106033 \\
        \hline
        Merge Sort & 0,000002 & 0,000009 & 0,000058 & 0,000551 & 0,006542 & 0,077718 \\
        \hline
    \end{tabular}
    \caption{Tempo de execução (em segundos) dos 5 algoritmos -- entrada Random}
    \label{tab:comparacao_random}
\end{table}


\qquad Para 1.000.000 de elementos, os destaques de cada tabela são:
na entrada crescente, o Bubble Sort (0,000321 s) e o Insertion Sort
(0,000584 s) foram os mais rápidos, por serem lineares nesse cenário,
seguidos pelo Shell Sort (0,008251 s) e pelo Merge Sort (0,021293 s),
enquanto o Selection Sort levou 934,9 s; na entrada decrescente, o Shell
Sort (0,011790 s) e o Merge Sort (0,020753 s) ficaram quatro ordens de
grandeza abaixo dos algoritmos quadráticos, entre os quais Bubble Sort
(155,6 s) e Insertion Sort (167,0 s) tiveram desempenho equivalente e o
Selection Sort foi o mais lento (937,3 s); na entrada randômica, o
Merge Sort foi o mais rápido (0,077718 s), seguido pelo Shell Sort
(0,106033 s), muito à frente do Insertion Sort (80,6 s), do Selection
Sort (932,5 s) e do Bubble Sort (938,7 s).
